% Minimal retained preamble for lane-p2 cards taken from
% sources/logic-fl02-conservation-of-ramsey-s-theorem-f/pinned-original-source.tex
% (Conservation of Ramsey's theorem for pairs, arXiv:2402.11616v3).
%
% Copied verbatim from the author's own file:
%   line 2      : \documentclass{article}
%   lines 16-25 : the author's theorem declarations
%   lines 46-88 : the author's macro definitions for the reverse-mathematics principles
%   line 91     : \card
%   lines 104-107: the natural sums and products, which rely on mathabx (source line 103)
%   line 116    : \qt
%
% Omitted author preamble lines and why (nothing was renamed or redefined):
%   line 3 (graphicx), 7 (url), 9 (cancel): no figure, url or cancel use in the retained spans;
%   line 6 (cleveref): the assembler adds its own cleveref machinery, a second load clashes;
%   line 8 (\usepackage[utf8x]{inputenc}): inputenc with the utf8x option fights the forced
%       xelatex engine;
%   lines 28-37 (the rep@theorem machinery and \newreptheorem{maintheorem}): the retained spans
%       use no rep-theorem environment;
%   lines 40-43 (xcolor and the three author comment macros): not used by the retained spans;
%   lines 93-99 (the LS2 font encoding block for \smashtimes based on stix): the host TeX Live has
%       no stix directory and the author himself never uses the result -- line 105 maps \ntimes to
%       \dottimes, and line 110 shows \smashtimes already commented out in the source.
% mathabx is retained because \dotplus, \dottimes and \dotdiv come from it.
\documentclass{article}
\usepackage{amsmath,amssymb,amsthm}
\usepackage{mathabx}
\usepackage{cancel}
\usepackage{cleveref}
\usepackage{url}

\newtheorem{theorem}{Theorem}
\numberwithin{theorem}{section}
\newtheorem{maintheorem}[theorem]{Main Theorem}
\newtheorem{lemma}[theorem]{Lemma}
\newtheorem{question}[theorem]{Question}
\newtheorem{proposition}[theorem]{Proposition}
\newtheorem{statement}[theorem]{Statement}
\newtheorem{remark}[theorem]{Remark}
\newtheorem{definition}[theorem]{Definition}
\newtheorem{corollary}[theorem]{Corollary}

\newcommand{\RCA}{\mathsf{RCA}}
\newcommand{\WKL}{\mathsf{WKL}}
\newcommand{\SRT}{\mathsf{SRT}}
\newcommand{\EM}{\mathsf{EM}}
\newcommand{\SEM}{\mathsf{SEM}}
\newcommand{\ADS}{\mathsf{ADS}}
\newcommand{\SCAC}{\mathsf{SCAC}}
\newcommand{\SADS}{\mathsf{SADS}}
\newcommand{\COH}{\mathsf{COH}}
\newcommand{\ACA}{\mathsf{ACA}}
\newcommand{\qvdash}{\operatorname{{?}{\vdash}}}
\newcommand{\nqvdash}{\operatorname{{?}{\nvdash}}}
\newcommand{\BSig}{\mathsf{B}\Sigma^0}
\newcommand{\ISig}{\mathsf{I}\Sigma^0}
\newcommand{\LSig}{\mathsf{L}\Sigma^0}
\newcommand{\BPi}{\mathsf{B}\Pi^0}
\newcommand{\IPi}{\mathsf{I}\Pi^0}
\newcommand{\LPi}{\mathsf{L}\Pi^0}
\newcommand{\IDel}{\mathsf{I}\Delta^0}
\newcommand{\BCDe}{\mathsf{BC}\Delta^0}
\newcommand{\WF}{\mathsf{WF}}
\newcommand{\RT}{\mathsf{RT}}
\newcommand{\BME}{\mathsf{BME}}
\newcommand{\TT}{\mathsf{TT}}
\renewcommand{\O}{\mathcal{O}}
\renewcommand{\L}{\mathcal{L}}
\newcommand{\C}{\mathcal{C}}
\newcommand{\R}{\mathcal{R}}
\renewcommand{\S}{\mathcal{S}}
\newcommand{\T}{\mathcal{T}}
\renewcommand{\P}{\mathcal{P}}
\newcommand{\M}{\mathcal{M}}
\newcommand{\Nc}{\mathcal{N}}
\renewcommand{\P}{\mathcal{P}}
\newcommand{\Q}{\mathcal{Q}}
\newcommand{\NN}{\mathbb{N}}
\newcommand{\PP}{\mathbb{P}}
\newcommand{\uh}[0]{{\upharpoonright}}
\newcommand{\card}{\operatorname{card}}
\newcommand{\nplus}{\dotplus}
\newcommand{\ntimes}{\dottimes}
\newcommand{\nminus}{\dotdiv}
\newcommand{\nsum}{\dot{\sum}}
\def\qt#1{``#1''}
\title{Ramsey's theorem for pairs preserves well-foundedness for the induction and collection schemes}
\author{Quentin Le Houérou, Ludovic Levy Patey, Keita Yokoyama}
\date{arxiv:2402.11616v3; CC BY 4.0}
\begin{document}
\maketitle
\noindent\textit{Editorial note.} The counted unit is the article's own Corollary cor:rt22-wf at source lines 1384-1391 together with the proof environment that begins on the next line (1392-1401); the pairing is the ordinary adjacent one.\par\medskip
\noindent\textit{Editorial note.} The retained context is the statement of Proposition prop:rt22-wf at source lines 1356-1363, which the counted proof applies; that proposition is counted by a sibling card of this lane, its proof is not retained, and the counted intervals are disjoint.\par\medskip
\noindent\textit{Editorial note.} The counted proof carries the author citation of Belanger inside its first line. The author's own bibliography file biblio.bib from the pinned e-print of arXiv:2402.11616v3 is attached to this task unedited and complete, with the bibtex backend, so that the citation resolves.\par\medskip
\noindent\textit{Editorial note.} The article has an earlier commented-out version of this same corollary at source lines 1372-1382; it is commented out by the author and is not author text for our purposes, and it is not retained.\par\medskip
\noindent\textit{Editorial note.} The retained preamble reproduces verbatim the author's document class at source line 2, his theorem declarations at source lines 16-25, his macro definitions at source lines 46-88 and 91, and his natural-operation macros at source lines 104-107, and it retains mathabx and cleveref; the stix-dependent LS2 block and the inputenc load with the utf8x option are not retained.\par\medskip
\noindent\textit{Editorial note.} The single bibliography entry needed by this card is provided as a task-directory fragment holding the entry exactly as the bibtex run had formatted it from the author own biblio.bib; the bibtex backend is not used because the harness order is xelatex twice and then bibtex, so a generated bbl is never read back and the citation would stay undefined in the final pass.\par\medskip
\noindent\textit{Allowed standard prerequisites (editorial).}\begin{itemize}
\item The RT22 proposition controlling the jump. For a countable model of RCA\_0 plus collection for Sigma\textasciicircum{}0\_2 plus well-foundedness of omega\_4 to the alpha, topped by a set Y, and every colouring and every set P above the jump of Y with the model of P satisfying RCA\_0-star, there is a homogeneous set G with P above the jump of G plus Y and with the model of G satisfying collection and well-foundedness of alpha; retained with this card at source lines 1356-1363.. The counted corollary obtains its auxiliary set P and then applies exactly this proposition; the corollary differs from it by replacing collection with an arbitrary induction or collection scheme at a finite level and by dropping the auxiliary set.
\item Omega-extension for finite levels of induction and collection. The proof invokes a result attributed to Harrington, Simpson-Smith and Hajek, recorded by Belanger, that every countable topped model of RCA\_0-star plus the scheme at level n minus one can be omega-extended into a model of the same theory plus weak Konig's lemma; the citation is carried in the proof and its entry is attached with this card.. It produces the set P above the jump of Y whose model satisfies the scheme at level n minus one, which is what the RT22 proposition needs.
\item Transfer of a finite-level scheme along an omega-extension. If a Sigma\textasciicircum{}0\_n formula holds in the extension of G and every such formula is Sigma\textasciicircum{}0\_\{n-1\} in the model of the auxiliary set P, then satisfaction of the scheme at level n minus one in that model gives satisfaction of the scheme at level n in the model of G; the proof states this transfer in its last two lines.. It is the step that raises the scheme from level n minus one to level n, which is the difference between the hypothesis and the conclusion of the counted corollary.
\end{itemize}\medskip
\begin{proposition}\label[proposition]{prop:rt22-wf}
Consider a countable model $\M = (M,S) \models \RCA_0 + \BSig_2 + \WF(\omega_4^{\alpha})$ topped by a set~$Y \in S$, where~$\alpha \in M$ is an ordinal~$\leq \epsilon_0$. Then, for every coloring~$f : [M]^2 \to 2$ in~$S$ and every set~$P \gg Y'$ such that $\M[P] \models \RCA_0^*$, there exists $G \subseteq M$ such that 
\begin{enumerate}
    \item $G$ is an $M$-infinite $f$-homogeneous set ;
    \item $P \geq_T (G \oplus Y)'$ ;
    \item $\M[G] \models \RCA_0 + \BSig_2 + \WF(\alpha)$.
\end{enumerate}
\end{proposition}

\begin{corollary}\label[corollary]{cor:rt22-wf}
Fix~$n \geq 2$. Let~$T_n$ be either~$\BSig_n$ or~$\ISig_n$.
Consider a countable model $\M = (M,S) \models \RCA_0 + T_n + \WF(\omega_4^{\alpha})$ topped by a set~$Y \in S$, where~$\alpha \in M$ is an ordinal~$\leq \epsilon_0$. Then, for every coloring~$f : [M]^2 \to 2$ in~$S$, there exists $G \subseteq M$ such that 
\begin{enumerate}
    \item $G$ is an $M$-infinite $f$-homogeneous set ;
    \item $\M[G] \models \RCA_0 + T_n + \WF(\alpha)$.
\end{enumerate}
\end{corollary}

\begin{proof}
By Harrington; Simpson–Smith; Hajek (see Belanger~\cite{Blanger2022ConservationTF}), every countable topped model of $\RCA_0^* + T_{n-1}$ can be $\omega$-extended into a model of~$\RCA_0^* + T_{n-1} + \WKL$. Since~$\M[Y'] \models \RCA_0^* + T_{n-1}$, there is some $P \subseteq M$ such that $P \gg Y'$ and $\M[P] \models \RCA_0^* + T_{n-1}$. By \Cref{prop:rt22-wf}, there exists $G \subseteq M$ such that 
\begin{enumerate}
    \item $G$ is an $M$-infinite $f$-homogeneous set ;
    \item $P \geq_T (G \oplus Y)'$ ;
    \item $\M[G] \models \RCA_0 + \BSig_2 + \WF(\alpha)$.
\end{enumerate}
Since every $\Sigma^0_n$ formula in~$\M[G]$ is $\Sigma^0_{n-1}$ in~$\M[P]$
and $\M[P] \models T_{n-1}$, $\M[G] \models T_n$.
\end{proof}

\begin{thebibliography}{99}
\bibitem{Blanger2022ConservationTF}
David~R. B{\'e}langer.
\newblock Conservation theorems for the cohesiveness principle.
\newblock 2022.
\newblock Available at \url{https://arxiv.org/abs/2212.13011}.
\end{thebibliography}
\end{document}
